% ===========================================================================
% Section 5: Main Results
%
% Robot-mode (Shapiro Step 3): state facts, do not argue; define concepts
% before use; no forward references to undefined constructs; plain precise
% language.
%
% Subsections drafted here (anchored by Tables 8, 9, 10):
%   5.1 First Stage
%   5.2 Population Effects
%   5.3 Sectoral Activity (renamed from "Employment and Sectoral Effects"
%       per PR #49 -- the paper-skeleton employment outcomes could not
%       be implemented; the industrial and agricultural census activity
%       outcomes are used instead. See table_10_sectoral.R header.)
%
% Subsections deferred (need more analysis before drafting):
%   5.4 Other Outcomes   (pending Table 11: education, migration,
%                        employment status)
%   5.5 Robustness       (pending Table 12: alternative theta, alternative
%                        hypo instruments, subsample stability)
%
% Numbers inline for now; moved to scalars.tex AutoFill macros once the
% analysis pipeline stabilizes.
%
% CITATION KEYS: none new beyond Sections 1-4.
% ===========================================================================

\section{Main Results}
\label{sec:results}

\subsection{First Stage}
\label{sec:first_stage}

Table~\ref{tab:first_stage} reports the first-stage regressions
corresponding to equation~\eqref{eq:first_stage}. Column~(1) uses only
the Larkin Plan instrument $\Delta \ln \mathrm{MA}^{\mathrm{LP}}_i$,
column~(2) uses only the hypothetical-road instrument
$\Delta \ln \mathrm{MA}^{\mathrm{hypo}}_i$, and column~(3) uses both
jointly. All columns include baseline log market access, baseline log
population, and the six standardized geographic controls. Standard
errors are heteroskedasticity-robust.

The Larkin Plan instrument has a coefficient of 0.240 with a robust
standard error of 0.054 in column~(1). The corresponding first-stage
$F$ statistic is 19.97. The hypothetical-road instrument has a
coefficient of 0.123 (0.069) in column~(2), with a first-stage $F$ of
3.14. In the joint specification in column~(3), both instruments
retain predictive power: the Larkin Plan coefficient is 0.254 (0.055)
and the hypothetical-road coefficient is 0.146 (0.063); the joint $F$
is 13.05.

The Larkin Plan instrument alone exceeds the Stock--Yogo rule-of-thumb
threshold of 10. The hypothetical-road instrument alone does not. The
joint specification exceeds the threshold and is adopted as the main
specification throughout this section. The hypothetical-road-only
column is reported so that the two sources of variation can be
inspected separately.

\subsection{Population Effects}
\label{sec:population}

Table~\ref{tab:population_iv} reports estimates of equation~\eqref{eq:main}
for four population outcomes: the log change in total population, in
urban population, and in rural population between 1960 and 1991, and
the level change in the urban share over the same period. Each
outcome panel reports four columns: OLS in column~(1), instrumental
variables using the Larkin Plan instrument in column~(2), the
hypothetical-road instrument in column~(3), and both instruments in
column~(4).

For total population (Panel~A), the OLS coefficient on $\Delta \ln
\mathrm{MA}^{\mathrm{full}}_i$ is 0.022 (0.012), significant at the
ten-percent level. The IV-LP estimate is 0.042 (0.042), the IV-Hypo
estimate is 0.059 (0.082), and the combined IV estimate is 0.046
(0.033). The sample is 309 districts.

For urban population (Panel~B), the OLS coefficient is 0.030 (0.013),
significant at the five-percent level. The IV-LP estimate is 0.062
(0.046), the IV-Hypo estimate is 0.049 (0.129), and the combined IV
estimate is 0.060 (0.040). The sample is 284 districts, reflecting
districts with positive urban population in both census years.

For rural population (Panel~C), the OLS coefficient is 0.031 (0.025).
The IV-LP estimate is 0.107 (0.108), the IV-Hypo estimate is 0.119
(0.165), and the combined IV estimate is 0.111 (0.087). The sample is
272 districts.

For the urban share (Panel~D), all four estimates are close to zero
and statistically indistinguishable from zero (OLS: 0.004 (0.005);
IV-Both: 0.002 (0.017)). A one-unit increase in $\Delta \ln
\mathrm{MA}^{\mathrm{full}}_i$ corresponds to a change in the urban
share of less than one percentage point.

Across the three log-change outcomes, the IV point estimates exceed
the OLS estimates by a factor of approximately two to four. This
pattern is consistent with a downward bias in OLS that would arise if
infrastructure changes were concentrated in districts with declining
pre-period trends; the direction of bias was a possibility raised in
Section~\ref{sec:identification}. The standard errors on the IV
estimates are large enough that individual point estimates are not
statistically distinguishable from zero at conventional levels. The
IV-Both column reports the narrowest confidence intervals among the
three IV columns.

% [PLACEHOLDER: A1 (OLS vs IV bias direction) — 1-2 paragraphs on the
%  direction and interpretation of the OLS-IV gap. Deferred until the
%  robustness in Section 5.5 is run.]

\subsection{Sectoral Activity}
\label{sec:sectoral}

Table~\ref{tab:sectoral_iv} reports estimates of
equation~\eqref{eq:main} for five sectoral outcomes drawn from the
industrial and agricultural censuses. Panel~A covers three
manufacturing outcomes from the 1954 and 1985 industrial censuses:
the log change in the number of establishments, in the value of
production, and in the total wage mass. Panel~B covers two
agricultural outcomes from the 1960 and 1988 agricultural censuses:
the log change in the number of farms and in the total farmed area.
Column structure matches Table~\ref{tab:population_iv}.

For manufacturing establishments, the OLS estimate is 0.025 (0.017)
and the combined IV estimate is 0.055 (0.061). Neither is
statistically distinguishable from zero.

For manufacturing production value, the OLS estimate is 0.060
(0.037). The IV-LP estimate is 0.239 (0.138), significant at the
ten-percent level. The combined IV estimate is 0.259 (0.119),
significant at the five-percent level. The sample is 308 districts.

For manufacturing wage mass, the OLS estimate is 0.092 (0.035),
significant at the one-percent level. The IV-LP estimate is 0.258
(0.146), significant at the ten-percent level. The combined IV
estimate is 0.330 (0.126), significant at the one-percent level.
The sample is 307 districts.

For the number of agricultural farms (Panel~B), the OLS estimate is
$-0.011$ (0.012). The combined IV estimate is 0.095 (0.079).
Neither is statistically distinguishable from zero.

For total agricultural farmed area, the OLS estimate is $-0.001$
(0.018). The combined IV estimate is 0.104 (0.078). Neither is
statistically distinguishable from zero.

The estimates differ substantially across the two panels. The
manufacturing outcomes related to activity intensity --- production
value and wage mass --- respond strongly and positively to
$\Delta \ln \mathrm{MA}^{\mathrm{full}}_i$: combined-IV point estimates
are 0.259 and 0.330, respectively. The number of manufacturing
establishments does not respond detectably. The agricultural
outcomes show point estimates close to zero with wide confidence
intervals. Section~\ref{sec:mechanisms} returns to the question of
whether this pattern is consistent with the scale-economies
hypothesis introduced in Section~\ref{sec:history}.

% [PLACEHOLDER: A2 (sectoral patterns and scale economies) — 2 paragraphs
%  connecting these empirical patterns to the Baumgartner-Palazzo cost
%  differences. Deferred until Section 7 interpretation is drafted,
%  because the scale-economies argument has to be articulated once,
%  consistently, across Sections 5, 6, and 7.]

% [PLACEHOLDER: Note on the OLS--IV gap for manufacturing wage mass.
%  The gap is large (0.092 OLS vs 0.330 IV-Both, a factor of 3.6).
%  One sentence on why: consistent with the downward-OLS-bias argument
%  from Section 5.2; specific to wage mass, possibly because wage
%  levels were rising in declining districts as worker composition
%  changed. Discuss with Cote before drafting.]

\subsection{Other Outcomes}
\label{sec:other_outcomes}

% [PLACEHOLDER: Section 5.4 pending. Depends on Table 11 (education,
%  migration, employment status from IPUMS). Draft after Table 11 runs.]

\subsection{Robustness}
\label{sec:results_robustness}

% [PLACEHOLDER: Section 5.5 pending. Depends on Table 12 (alternative
%  theta = 8.11, alternative hypo instruments, subsample stability).
%  Draft after Table 12 runs.]
